% =============================================================================
% UBS COE XD - SPEED-FORMAT PROPOSAL
% =============================================================================
\documentclass[11pt,a4paper]{article}

\usepackage[T1]{fontenc}
\usepackage[utf8]{inputenc}
\usepackage[ngerman]{babel}
\usepackage[top=25mm,bottom=20mm,left=25mm,right=20mm]{geometry}
\usepackage{xcolor}
\usepackage{helvet}
\usepackage{tcolorbox}
\usepackage{booktabs}
\usepackage{colortbl}
\usepackage{hyperref}
\usepackage{titlesec}
\usepackage{fancyhdr}
\usepackage{lastpage}
\usepackage{enumitem}
\usepackage{array}
\usepackage{amssymb}
\usepackage{tikz}

% Colors
\definecolor{fadarkblue}{RGB}{2,64,121}
\definecolor{falightblue}{RGB}{84,158,222}
\definecolor{fadarkgray}{RGB}{37,33,42}
\definecolor{falightgray}{RGB}{243,245,247}
\definecolor{famint}{RGB}{126,189,172}
\definecolor{faorange}{RGB}{222,157,62}
\definecolor{fawhite}{RGB}{255,255,255}

\renewcommand{\familydefault}{\sfdefault}
\AtBeginDocument{\color{fadarkgray}}

\titleformat{\section}{\bfseries\Large\color{fadarkblue}}{\thesection}{1em}{}
\titleformat{\subsection}{\bfseries\large\color{fadarkblue}}{\thesubsection}{1em}{}

\hypersetup{colorlinks=true,linkcolor=fadarkblue,urlcolor=fadarkblue}

\pagestyle{fancy}
\fancyhf{}
\fancyhead[L]{\small\color{fadarkgray}Speed-Format | CoE Experience Design}
\fancyhead[R]{\small\color{fadarkgray}Januar 2026}
\fancyfoot[L]{\small\color{fadarkblue}FehrAdvice \& Partners AG}
\fancyfoot[R]{\small\color{fadarkgray}Seite \thepage\ von \pageref{LastPage}}
\renewcommand{\headrulewidth}{0.5pt}

\newtcolorbox{fainfobox}[1][]{colback=falightgray,colframe=fadarkblue,coltitle=fawhite,fonttitle=\bfseries,title=#1}
\newtcolorbox{fasuccessbox}[1][]{colback=famint!10,colframe=famint,coltitle=fawhite,fonttitle=\bfseries,title=#1}
\newcommand{\fahead}[1]{\textcolor{fawhite}{\bfseries #1}}

\begin{document}

% =============================================================================
% TITLE PAGE
% =============================================================================
\begin{titlepage}
\begin{center}
\vspace*{1.5cm}

{\Huge\color{fadarkblue}\textbf{FehrAdvice}}\\[0.2cm]
{\large\color{falightblue}Behavioral Economics}

\vspace{1.5cm}

{\Huge\bfseries\color{fadarkblue}Das Speed-Format\par}
\vspace{0.5cm}
{\Large\color{falightblue}Verhaltensökonomische Antworten in 2--5 Tagen\par}

\vspace{0.8cm}
{\color{fadarkblue}\rule{0.8\textwidth}{2pt}}
\vspace{0.8cm}

{\large für das\\[0.3cm]\textbf{Center of Excellence Experience Design}\\[0.3cm]UBS Switzerland\par}

\vfill

{\small\color{fadarkgray}
FehrAdvice \& Partners AG | Januar 2026\\[0.3cm]
\textit{Vertraulich}
}

\end{center}
\end{titlepage}

% =============================================================================
\section{Das Angebot auf einen Blick}
% =============================================================================

\begin{fainfobox}[Was ist das Speed-Format?]
\textbf{Verhaltensökonomische Analysen für das gesamte CoE XD Team -- Antworten in 2--5 Tagen.}

\vspace{0.3cm}
Jede:r Mitarbeitende im CoE kann eine Fragestellung einreichen. FehrAdvice liefert innerhalb von 2--5 Tagen einen strukturierten Report mit verhaltensökonomischer Analyse und konkreten Empfehlungen.
\end{fainfobox}

\vspace{0.5cm}

\begin{center}
\begin{tabular}{p{4cm}p{9cm}}
\toprule
\rowcolor{fadarkblue}
\fahead{Aspekt} & \fahead{Speed-Format} \\
\midrule
\rowcolor{fawhite}
\textbf{Wer kann fragen?} & Alle Mitarbeitenden des CoE XD \\
\rowcolor{falightgray}
\textbf{Turnaround} & 2--5 Arbeitstage \\
\rowcolor{fawhite}
\textbf{Output} & Strukturierter Report (5--10 Seiten) \\
\rowcolor{falightgray}
\textbf{Interaktion} & 2 synchrone Sessions + asynchrone Arbeit \\
\bottomrule
\end{tabular}
\end{center}

\vspace{0.8cm}

\subsection{Das ist es -- Das ist es nicht}

\begin{center}
\begin{tabular}{p{6.5cm}p{6.5cm}}
\rowcolor{famint!20}
\textbf{\color{famint}$\checkmark$ Das ist es} & \textbf{\color{faorange}$\times$ Das ist es nicht} \\
\midrule
Definierter Rahmen & Grossprojekt \\
Klare Lieferobjekte & Ad-hoc Anfragen ohne Struktur \\
Fester Zeitrahmen (2--5 Tage) & Open-End Beratung \\
Scope Freeze nach Briefing & Laufende Scope-Erweiterung \\
\end{tabular}
\end{center}

\newpage
% =============================================================================
\section{Der Prozess: 5 Phasen in 2--5 Tagen}
% =============================================================================

\begin{center}
\begin{tabular}{clp{8cm}}
\toprule
\rowcolor{fadarkblue}
\fahead{Phase} & \fahead{Was} & \fahead{Beschreibung} \\
\midrule
\rowcolor{fawhite}
\textbf{1} & \textbf{Briefing-Call} (30 min) & Synchroner Auftakt: Fragestellung schärfen, Scope definieren, Kontext klären \\
\rowcolor{falightgray}
\textbf{2} & \textbf{Vorstrukturierung} & FehrAdvice entwickelt asynchron einen ersten Strukturvorschlag \\
\rowcolor{fawhite}
\textbf{3} & \textbf{Session 1} (60 min) & Abstimmung des Vorschlags und der finalen Lieferobjekte \\
\rowcolor{falightgray}
\textbf{4} & \textbf{Detailausarbeitung} (1--3 Tage) & Asynchrone Ausarbeitung durch FehrAdvice \\
\rowcolor{fawhite}
\textbf{5} & \textbf{Session 2} (60 min) & Präsentation, Validierung, finale Freigabe \\
\bottomrule
\end{tabular}
\end{center}

\vspace{0.8cm}

\subsection{Leitplanken für den Erfolg}

\begin{itemize}[leftmargin=*]
    \item \textbf{Scope Freeze nach Phase 1:} Nach dem Briefing-Call werden keine neuen Fragestellungen mehr aufgenommen
    \item \textbf{Fester Zeitrahmen:} Das Engagement hat ein definiertes Ende -- keine Endlosberatung
    \item \textbf{Sessions zur Präzisierung:} Die Meetings dienen der Verfeinerung, nicht der kompletten Neuausrichtung
\end{itemize}

\vspace{0.5cm}

\begin{fasuccessbox}[Typische Fragestellungen]
\begin{itemize}[leftmargin=*]
    \item «Warum brechen Kund:innen an Punkt X in der Journey ab?»
    \item «Welche Defaults sollten wir im neuen Feature setzen?»
    \item «Wie reduzieren wir die Komplexität bei der Produktauswahl?»
    \item «Welche BE-Hebel können wir für höhere Feature-Adoption nutzen?»
\end{itemize}
\end{fasuccessbox}

\newpage
% =============================================================================
\section{Investition}
% =============================================================================

\begin{fainfobox}[Das Modell]
\begin{center}
{\LARGE\textbf{CHF 12'000.- / Monat}}\\[0.5cm]
{\large für bis zu \textbf{3 Fragestellungen} pro Monat}
\end{center}
\end{fainfobox}

\vspace{0.8cm}

\subsection{Was ist enthalten?}

\begin{itemize}[leftmargin=*]
    \item[$\checkmark$] Bis zu 3 Speed-Analysen pro Monat
    \item[$\checkmark$] Pro Analyse: Briefing (30 min) + 2 Sessions (je 60 min)
    \item[$\checkmark$] Verhaltensökonomische Analyse mit EBF-Framework
    \item[$\checkmark$] Strukturierter Report (5--10 Seiten) pro Fragestellung
    \item[$\checkmark$] Turnaround: 2--5 Arbeitstage pro Analyse
\end{itemize}

\vspace{0.5cm}

\subsection{Out of Scope}

\begin{itemize}[leftmargin=*]
    \item[$\times$] Empirische Primärerhebung (Umfragen, Experimente)
    \item[$\times$] Operative Umsetzung oder Design-Implementierung
    \item[$\times$] Mehr als 3 Fragestellungen pro Monat
\end{itemize}

\vspace{1cm}

% =============================================================================
\section{Nächste Schritte}
% =============================================================================

\begin{center}
\begin{tabular}{clp{8cm}}
\toprule
\rowcolor{fadarkblue}
\fahead{\#} & \fahead{Schritt} & \fahead{Details} \\
\midrule
\rowcolor{fawhite}
1 & Vereinbarung unterzeichnen & Start per sofort möglich \\
\rowcolor{falightgray}
2 & Kick-off mit dem CoE XD & Einführung für das Team (30 min) \\
\rowcolor{fawhite}
3 & Erste Frage einreichen & Start des ersten Speed-Zyklus \\
\bottomrule
\end{tabular}
\end{center}

\vfill

\begin{center}
\textcolor{fadarkblue}{\rule{0.8\textwidth}{1pt}}\\[0.5cm]
\textbf{FehrAdvice \& Partners AG}\\
Klausstrasse 4 | 8008 Zürich | www.fehradvice.com
\end{center}

\end{document}
